\documentclass[11pt,a4paper]{article}
\usepackage[T1]{fontenc}
\usepackage[utf8]{inputenc}
\usepackage{lmodern,microtype}
\usepackage[margin=25mm]{geometry}
\usepackage{amsmath,amssymb,amsthm,mathtools}
\usepackage{booktabs,longtable,array,enumitem,xcolor,listings,hyperref}
\hypersetup{colorlinks=true,linkcolor=blue!45!black,citecolor=blue!45!black,
 urlcolor=blue!45!black,pdftitle={Smaller Diophantine formulations of Con(ZFC)},
 pdfauthor={Research synthesis prepared for Vladimir Reshetnikov and ProveIt}}
\setlength{\emergencystretch}{3em}
\setlength{\parskip}{3pt}
\setlist{itemsep=2pt,topsep=4pt}
\newtheorem{theorem}{Theorem}[section]
\newtheorem{lemma}[theorem]{Lemma}
\newtheorem{proposition}[theorem]{Proposition}
\newtheorem{corollary}[theorem]{Corollary}
\theoremstyle{remark}\newtheorem{remark}[theorem]{Remark}
\newcommand{\Z}{\mathbb Z}
\newcommand{\N}{\mathbb N}
\newcommand{\Pos}{\mathbb Z_{>0}}
\newcommand{\Con}{\operatorname{Con}}
\newcommand{\pc}{\operatorname{pc}}
\newcommand{\bits}{\operatorname{bits}}
\newcommand{\band}{\mathbin{\&}}
\newcommand{\file}[1]{\nolinkurl{#1}}
\newcommand{\Msrc}{M_F^{\mathrm{src}}}
\newcommand{\Mnat}{M_F^{\mathrm{nat}}}
\newcommand{\RR}{\mathcal R}
\newcommand{\Arat}{\mathcal A}
\newcommand{\Brat}{\mathcal B}
\newcommand{\Drat}{\mathcal D}
\lstset{basicstyle=\ttfamily\footnotesize,breaklines=true,columns=fullflexible,
 keepspaces=true,frame=single}
\title{\textbf{Smaller Diophantine formulations of \(\Con(\mathrm{ZFC})\)}\\[5pt]
 \large Four-cell coefficients, a six-equation degree-28 system,\\
 and single equations in fourteen positive unknowns}
\author{Research synthesis prepared for Vladimir Reshetnikov\\
 and the ProveIt repository}
\date{22 October 2026\\[3pt]\small
 Repository snapshot \texttt{3b201cdea32d2c816f5f3d20734051f25fd3ef14}}
\begin{document}
\maketitle
\begin{abstract}
We combine four compatible reductions from the eight existing ProveIt
packages on Diophantine forms of ZFC consistency. The resulting family
improves several of their measured bounds. First, the input-free
construction and the ordinary strong Pell norm give six equations in
fifteen positive unknowns, of maximum degree 28, with a 61-operation
evaluator. The previous degree-28 system used sixteen unknowns, seven
equations and 63 operations. Its product finalizer has degree 99 and
69 operations; its sum of squares has degree 56 and 78 operations.
Second, transferring the existing quadratic-root elimination to the
input-free construction gives a single equation in fourteen positive
unknowns: either 73 operations and degree 342, or 75 operations and
degree 293. Third, the ordinary strong norm removes a register from the
existing quadratic lift, giving 24 quadratic equations in 33 positive
unknowns and a quartic sum of squares.

A separate compiler reduction replaces nine-cell windows by four directed
cells. Moving the predecessor initialization test to its predecessor
removes that coordinate as well as all four unused corners. The finite
tile alphabet decreases from about \(1.42\cdot10^{21}\) to 71,356,645.
The coefficient \(B-1\) now has exactly \(5^{28}\) binary digits instead
of \(5^{67}\), a reduction by \(5^{39}\). All coefficients are defined
by finite, unconditional recipes; their enormous expansions are not
materialized. The report supplies ordinary proofs, explicit arithmetic
circuits, the exact machine table, and reproducible standard-library
checks. These are improvements within the inspected constructions,
not global optimality theorems. The new arithmetic and compiler bridge
are not Lean/Rocq formalizations and do not prove consistency.
\end{abstract}
\clearpage\tableofcontents\newpage

\section{The result and the comparison being made}
Let \(\operatorname{Proof}_{\mathrm{ZFC}}(p,\bot)\) be the usual effective
finite-refutation predicate, in the presentation calibrated to K215.
As usual, \(\Con(\mathrm{ZFC})\) means
\(\neg\exists p\,\operatorname{Proof}_{\mathrm{ZFC}}(p,\bot)\).
The repository's endpoint identifies this refutation predicate with
halting of the literal blank-tape machine K215~\cite{endpoint}.

\begin{theorem}[Fixed consistency representations]\label{thm:main}
Fix the four integer coefficients of Section~\ref{sec:coefficients}.
Let \(\mathcal S_{15}\), \(P_{69}\), \(\Phi_{14}\), \(\Psi_{14}\),
and \(Q_{33}\) be the systems and polynomials defined below. ZFC proves
that each of the following assertions is equivalent to
\(\exists p\,\operatorname{Proof}_{\mathrm{ZFC}}(p,\bot)\):
\begin{enumerate}
\item \(\mathcal S_{15}\) has a solution in \(\Pos^{15}\);
\item \(P_{69}=0\) has a solution in \(\Pos^{15}\);
\item \(\Phi_{14}=0\), or equivalently \(\Psi_{14}=0\), has a solution
in \(\Pos^{14}\);
\item \(Q_{33}=0\) has a solution in \(\Pos^{33}\).
\end{enumerate}
Consequently, the absence of solutions of each displayed construction is
provably equivalent in ZFC to \(\Con(\mathrm{ZFC})\).
The explicitly defined nonnegative and unrestricted-integer variants in
Section~\ref{sec:domains} have the same property.
\end{theorem}

The proof is the composition of the ordinary K215 calibration, the
native Pell/mask compiler, and the changes proved here. The assertion
about \emph{provability in ZFC} uses the formalizability of those ordinary
arguments. It does not infer an internal ZFC theorem from the soundness
of Lean or Rocq. Section~\ref{sec:status} states the verification boundary.

\subsection{What ``simpler'' means here}
There are different measures: number of unknowns, number of equations,
degree, arithmetic operations, coefficient height, and proof complexity.
The operation convention follows the existing packages: each binary
\(+,-,\times\), including multiplication by a fixed numeral, costs one;
reading a constant, register reuse, and equality comparison cost zero.
Intermediate expressions are substituted polynomial abbreviations, not
additional existential variables. A sum-of-squares finalizer is fully
charged. Coefficient bit length is reported separately rather than hidden
inside the operation figure.

\begin{center}\small
\begin{tabular}{lrrrr}
\toprule
Construction & Positive unknowns & Equations & Degree & Operations\\
\midrule
Existing 02 degree-28 system & 16 & 7 & 28 & 63\\
Existing 06/08 basic system & 15 & 6 & 60 & 61\\
\textbf{New \(\mathcal S_{15}\)} & \textbf{15} & \textbf{6} & \textbf{28} & \textbf{61}\\
Existing 06/08 product & 15 & 1 & 155 & 67\\
Existing 06 reduced product & 15 & 1 & 123 & 68\\
New \(P_{69}\) & 15 & 1 & 99 & 69\\
Existing 02 sum of squares & 16 & 1 & 56 & 83\\
New sum of squares of \(\mathcal S_{15}\) & 15 & 1 & 56 & 78\\
New \(\Psi_{14}\) & 14 & 1 & 342 & 73\\
New \(\Phi_{14}\) & 14 & 1 & 293 & 75\\
Existing 08 quadratic lift & 34 & 25 & 2 & --\\
New quadratic lift & 33 & 24 & 2 & --\\
New \(Q_{33}\), its sum of squares & 33 & 1 & 4 & --\\
\bottomrule
\end{tabular}
\end{center}
The systems' degree entries are maximum residual degrees; the single
polynomials' entries are exact. The original 67-operation polynomial is
still the shortest operation count among these single-polynomial
evaluators, and can also use the new coefficients. Fourteen is the
smallest positive-variable count among the eight inspected packages and
these additions, but it is not a global record: classical nine-unknown
representations and the repository's existence theorem are a different
comparison class~\cite{reconcile,nine}. No minimality assertion is made.

\subsection{How the eight packages contribute}
The baseline reconciliation is taken from the pinned repository, then
the relevant formulas and proof dependencies were read directly. Package
01 gives conservative fixed-input and quadratic alternatives. Package
02 supplies the ordinary strong norm and its integrality/rank proof.
Package 03 combines fixed markers with root elimination. Package 04
explores direct two-sided normalization and an affine marker guard.
Package 05 provides particularly explicit root-elimination formulas and
integer conversions. Packages 06 and 08 independently remove the input
marker and reach the same 67-operation arithmetic skeleton; 06 has the
smaller, fully numerical raw-K215 coefficient recipe. Package 07 gives
alternative scales and degree/operation tradeoffs. Package 08 supplies
the 34-variable quadratic lift.

The synthesis is therefore specific: use 06's input-free interface,
02's ordinary norm, 05's root elimination where variable count matters,
and 08's quadratic organization. The four-cell coefficient construction
is an additional reduction. Coefficients from separate compiler recipes
must never be mixed. All four ports in this report come from the single
new recipe in Section~\ref{sec:coefficients}; the complete package-06
tuple remains an independently valid legacy choice for the algebra.

\section{The recommended six-equation system}\label{sec:system}
The fifteen unknowns are
\[
 (J,F,\alpha,z,f,h,i,T,s,w,\tau,\eta,\zeta,y,\gamma)\in\Pos^{15}.
\]
The fixed coefficients are \(\beta,K,M_C,\Msrc\). Define
\begin{equation}
\begin{gathered}
 q=\beta J+1,\quad C=q-F-\alpha,\quad X=wq,\quad Y=sq^3,\\
 E=XY,\quad k=\eta+\zeta,\quad c=kY+\eta,\quad a=E+Y,\\
 A=a+2,\quad H=4a+3,\quad \Delta=a^2+H=A^2-1,\\
 D=X+ac+\gamma H,\quad G=(q-1)(q-F)+\alpha,\\
 R=G(q^2-1)+(M_C+q\Msrc)J,\quad L_0=EkY,\\
 V=c(Tf-1)-Rf^2,\quad U=ic^2,\quad W=\Delta(f^2-1).
\end{gathered}
\label{eq:definitions}
\end{equation}
The symbol \(W\) is only an auxiliary polynomial in this section; there
is no numerical-input marker. The six equations are
\begin{equation}
\boxed{\begin{aligned}
 \tau^2-L_0(L_0+k)&=1,\\
 D^2-\Delta c^2&=1,\\
 W(V^2-y^2)+y^2&=1,\\
 k-hE-R&=1,\\
 (K+w)C+q-F-z(q-1)&=1,\\
 U^2&=W.
\end{aligned}}\label{eq:system}
\end{equation}
Every exponent in these polynomials is a fixed integer. All equalities
are ordinary polynomial equalities over the stated domain.

Let \(N_0,N_1,N_a,N_k,N_t\) denote the first five left sides and put
\(N_o=1+U^2-W\). Then a single-polynomial alternative is
\begin{equation}
 \boxed{P_{69}=N_0N_1N_aN_kN_tN_o-1=0.}\label{eq:p69}
\end{equation}
Its positive zero set is exactly the solution set of~\eqref{eq:system},
not merely equivalent in nonemptiness. The sum of the six squared
residuals is another single equation, of degree 56.

\subsection{The normalized parent and the changed proof order}
For comparison, the input-free normalized parent has
\[
 S=(i\Delta c^2)^2,\quad N_a^{\rm n}=S(V^2-y^2)+y^2,
 \quad N_s^{\rm n}=\Delta f^2-S,
\]
and polynomial
\begin{equation}
 P_{67}^{\rm n}=N_0N_1N_a^{\rm n}N_kN_tN_s^{\rm n}-\Delta.
 \label{eq:normalized}
\end{equation}
Its strong equation is \(f^2-\Delta i^2c^4=1\). In contrast,
the ordinary equation in~\eqref{eq:system} is
\((ic^2)^2=\Delta(f^2-1)\). These two individual equations are not
unconditionally equivalent. The ordinary equation can have extra
rational Pell solutions when isolated. It becomes sufficient here
because the main Pell ordinate is already large. The central proof
change is to establish rank \emph{before} restoring the auxiliary
quotient. This is why the combination requires a proof rather than a
formal substitution.

\section{Proof of the ordinary-norm combination}\label{sec:ordinaryproof}
Write \(\chi_A(n)+\psi_A(n)\sqrt{A^2-1}=(A+\sqrt{A^2-1})^n\).
Positive Pell solutions at parameter \(A\ge2\) are classified by these
sequences. The two sequences satisfy the same second-order recurrence,
and \(\psi_A\) has the strong divisibility property.

\subsection{Unit signs and the boundary \texorpdfstring{\(C=0\)}{C=0}}
At a positive zero of~\eqref{eq:p69}, every factor is an integer unit.
The negative consecutive-product norm
\(x^2-u(u+1)v^2=-1\), for \(u\ge2\), is impossible: a solution
minimal in \(v>0\) is reduced by multiplication by
\(2u+1-2\sqrt{u(u+1)}\), whose new second coordinate is
\(0<(2u+1)v-2|x|<v\). Thus \(N_0=1\).
Since \(\Delta\equiv0\) or \(3\pmod4\), the main norm cannot
equal \(-1\). The ordinary strong factor also cannot be \(-1\):
when \(\Delta\equiv0\), it is \(1+U^2\pmod4\); when
\(\Delta\equiv3\), it is \(U^2+f^2\pmod4\). Neither is 3.
Hence \(N_o=1\) and \(W=U^2\). The auxiliary norm is consequently
congruent modulo four to either \(y^2\) or \(V^2\), excluding its
negative unit. We have
\[
 N_0=N_1=N_a=N_o=1,\qquad N_k=\epsilon,\quad N_t=\nu,
 \qquad\epsilon\nu=1,
 \quad\epsilon,\nu\in\{1,-1\}.
\]

Transport gives \(C\ge0\), because
\[
 q-F-z(q-1)=1-F-(z-1)(q-1)\le0
\]
and \(C\le-1\) would make \(N_t\le-(K+w)\le-2\).
One must not conclude \(C>0\) yet: the boundary
\(C=0,F=2,z=1,\alpha=q-2\) can make the transport unit negative.
The source-mask bounds now give
\begin{gather}
 1\le F\le q-1,\qquad q\le G\le q^2-q,\notag\\
 q(q^2-1)<R<q^4-q^3+q^2-1<q^4-2,\label{eq:outerbounds}\\
 R-2>3q+1,\qquad E\ge q^4>R+2,\qquad a>R+2.
\notag
\end{gather}
Here the shifted mask satisfies \(\Msrc<2(B-1)\), rather than
the unshifted bound \(\Msrc<B-1\). These weaker, correct input-free
bounds are sufficient throughout.

\subsection{Main-index growth independent of the auxiliary sign}
Put \(P_0=2XY^2+1\). The first and main norms give positive indices
\(n,p\) with
\[
 \tau=\chi_{P_0}(n),\quad k=2\psi_{P_0}(n),\qquad
 D=\chi_A(p),\quad c=\psi_A(p).
\]
Since \(P_0\equiv1\pmod E\), the index unit gives
\(2n\equiv R+\epsilon\pmod E\). The positive representative
\(R+\epsilon\) lies below \(E\), so
\(n\ge(R-1)/2\ge24\). Moreover \(P_0>A\) and \(c>k\)
force \(p>n\). Elementary Pell growth, for example
\(\psi_A(p)\ge(2A-1)^{p-1}\), gives
\begin{equation}
 c>A\Delta^2,\qquad c>2p,\qquad c>2R.\label{eq:growth}
\end{equation}
The last inequality also follows from \(k=R+\epsilon+hE\),
\(E>R+2\), and \(c>kY\). No sign of \(V\) was used.

\begin{lemma}[Ordinary strong rank]\label{lem:ordinaryrank}
If \(c=\psi_A(p)>A\Delta^2\) and
\((ic^2)^2=\Delta(f^2-1)\) with positive integers, then some \(m>0\)
satisfies
\[
 f=\chi_A(m),\qquad p\mid m,\qquad c\mid m,
 \qquad ic^2=\Delta\psi_A(m).
\]
In particular \(c^2\mid f^2-1\).
\end{lemma}
\begin{proof}
This is the ordinary-rank lemma of package 02 and its elementary source
proof~\cite{ordinary,relaxed}; we include the argument needed here.
Write \(\Delta=d_0s_0^2\) with \(d_0>1\) squarefree. Let
\(F_0+y_0\sqrt{d_0}\) be the least positive integer-coefficient
norm-one unit. All positive integer-coefficient units in that ring are
its powers, by division by an appropriate power of the least unit.
For some \(e\ge1\), with \(S_0=\psi_{F_0}(e)\),
\[
 A=\chi_{F_0}(e),\quad\Delta=(F_0^2-1)S_0^2,
 \qquad\psi_{F_0}(ep)=S_0c.
\]
The ordinary square equation supplies another unit and an index \(v>0\):
\[
 f=\chi_{F_0}(v),\qquad ic^2=T_0\psi_{F_0}(v),
 \qquad T_0=(F_0^2-1)S_0\le\Delta.
\]
To justify its integrality, \(s_0^2\mid(ic^2)^2\) implies
\(s_0\mid ic^2\); squarefreeness then implies
\(d_0\mid ic^2/s_0\), giving the integer second coordinate.

Let \(g=\gcd(v,ep)\) and \(h_0=\psi_{F_0}(g)\).
Strong divisibility gives
\[
 c\gcd(c,\Delta)=\gcd(c^2,T_0\psi_{F_0}(ep))\mid T_0h_0,
 \qquad h_0\ge c/\Delta.
\]
If \(g<ep\), it is a proper divisor and \(2g\le ep\); Pell
duplication gives \(2h_0^2<\psi_{F_0}(2g)\le S_0c<Ac\).
This contradicts \(c>A\Delta^2\) and \(h_0\ge c/\Delta\).
Thus \(ep\mid v\); write \(v=em\). We obtain \(p\mid m\),
\(f=\chi_A(m)\), and \(ic^2=\Delta\psi_A(m)\).

Write \(m=pt\). Expansion of
\((\chi_A(p)+c\sqrt\Delta)^t\) gives
\[
 \frac{\psi_A(pt)}c\equiv t\chi_A(p)^{t-1}\pmod c.
\]
Since \(c^2\mid\Delta\psi_A(m)\) and \(\gcd(c,\chi_A(p))=1\),
we have \(c\mid\Delta t\). The Pell expansion modulo \(\Delta\)
gives \(\gcd(c,\Delta)=\gcd(p,\Delta)\), which divides \(p\).
Consequently \(c\mid pt=m\). Finally \(p\mid m\) implies
\(c\mid\psi_A(m)\), and the Pell norm implies
\(c^2\mid\Delta\psi_A(m)^2=f^2-1\).
\end{proof}

\subsection{Restoring the deleted coordinates}
The ordinary equation and~\eqref{eq:growth} imply \(f>2c\) and
\(f^2>\Delta+c\). Since \(R<a<\Delta\),
\begin{equation}
 W-(Rf^2+c)=(\Delta-R)f^2-\Delta-c>0.\label{eq:gap}
\end{equation}
The auxiliary norm is \(WV^2-(W-1)y^2=1\). It excludes \(V=0\).
For \(v=|V|>1\), the identity
\(v^2-1=(W-1)(y^2-v^2)\) and \(y\ge v+1\) imply
\(v\ge2W-1\). But \(T>0\) and~\eqref{eq:gap} give
\(V>-(Rf^2+c)>-W\), excluding negative \(V\) of magnitude greater
than one. The congruence \(V\equiv-c\pmod f\), with \(f>2c>2\),
excludes \(V=\pm1\). Thus \(V>1\).

Define, in the proof only,
\[
 o=cT-Rf,\qquad j=Tf-1-R(f^2-1)/c.
\]
Lemma~\ref{lem:ordinaryrank} proves that \(j\) is integral. The identities
\(of=V+c\) and \(jc=V+R\) prove \(o,j>0\), and
\(V=of-c=jc-R\). Division has not been introduced into the equation.

The standard two-congruence auxiliary lemma now gives \(p=R\).
For clarity, it classifies the auxiliary norm at base \(U\): the index
is odd, say \(\ell=2v+1\), and
\(V=Q_v(U^2)\), where
\[
 Q_v(1-A^2)=(-1)^v\psi_A(2v+1),\qquad
 Q_v(0)=(-1)^v(2v+1).
\]
Reducing modulo \(f=\chi_A(m)\), then using Pell step-down, gives
\(\ell\equiv\pm p\pmod m\). Reducing modulo \(c\), where
\(c\mid U,m\), gives \(R\equiv\pm p\pmod c\). The bounds
\(0<R,p<c/2\) force \(p=R\). This is exactly the inherited
interface~\cite{ordinary}; all its positivity and integrality premises
have now been supplied.

Write \(2n=R+\epsilon+vE\), \(v\ge0\). If \(v\ge1\),
\(E>R+2\) contradicts \(n<p=R\). Hence \(2n=R+\epsilon\).
The branch \(\epsilon=-1\) gives \(p=2n+1\). Since
\(2A^2-1>P_0\) and \(A>Y+1\), duplication gives
\(\psi_A(2n)\ge Ak\), contradicting
\(c=\psi_A(2n+1)<k(Y+1)\). Therefore \(\epsilon=\nu=1\).
Only now does transport exclude \(C=0\) and prove \(C>0\).

\subsection{Power recovery, masks, and completeness}
The retained Pell ratio argument recovers \(X=2^R\); hence
\(q=B^N\) and \(J=(q-1)/(B-1)\). Exact half-binomial extraction
gives \(\pc(R)\ge3dN+2\). Define
\[
 S_w=(C-1)+qF,\qquad
 T_w=(M_CJ+1)+q(\Mnat J-1).
\]
The source's shifted port gives the exact identity
\begin{equation}
 R=(q^2-S_w)(q^2-1)+T_w.\label{eq:wordidentity}
\end{equation}
Since \(0<S_w<q^2\), the inverse-population lemma forces
\begin{equation}
 (C-1)\band(M_CJ+1)=0,
 \qquad F\band(\Mnat J-1)=0.\label{eq:masksrecovered}
\end{equation}
The native compiler then decodes a periodic marked computation; the
new coefficient proof is supplied in Section~\ref{sec:coefficients}.
It yields a finite blank halting run of K215.

Conversely, a finite halting run gives a normalized positive zero of
\eqref{eq:normalized} by the native compiler. Keep every coordinate
except replace \(i_{\rm n}\) by \(i_{\rm o}=\Delta i_{\rm n}\).
Then \((i_{\rm o}c^2)^2=\Delta(f^2-1)\) and the auxiliary
coefficient is unchanged. This gives a positive solution of
\eqref{eq:system}. Thus the ordinary and normalized systems have the
same existence condition, although their individual \(i\)-coordinates
need not agree. The unit argument proves the asserted exact equality
of the ordinary system and product zero sets.

\subsection{Exact degrees and operation counts}
Write \(q_0=\beta J\), \(k_0=\eta+\zeta\), and
\(T_0=w(q_0-F-\alpha)-zq_0\) for leading forms only.
The main norm has the all-value cancellation
\[
 (X+ac+\gamma H)^2-(a^2+H)c^2
 =X^2+2X\gamma H+(\gamma H)^2+2acX+2ac\gamma H-Hc^2.
\]
The six residual degrees are \(22,18,28,7,2,22\).
The product's degree-99 leading form is
\begin{equation}
 8q_0^{55}h\gamma i^2k_0^9w^7s^{16}T^2f^4T_0.\label{eq:p69leader}
\end{equation}
For example the monomial
\(J^{56}h\gamma i^2\eta^9w^7s^{16}T^2f^4z\)
has coefficient \(-8\beta^{56}\ne0\). The unique degree-28
residual supplies a nonzero degree-56 square, proving the SOS degree.
The emitted source ledgers are
\[
 \mathcal S_{15}:34M+27A=61,\quad
 P_{69}:39M+30A=69,\quad
 \mathrm{SOS}:40M+38A=78.
\]
The files give every gate and certify that every gate and supplied port
is live. Diagnostic line specializations supplement these uniform
degree proofs; they do not replace them.

\section{Single equations in fourteen positive unknowns}\label{sec:root}
Return to the normalized parent~\eqref{eq:normalized} and delete \(f\)
from its fifteen-variable list. All definitions independent of \(f\)
are retained. Put
\begin{align}
 Q&=(\Delta ic^2)^2,&\Drat&=1+\Delta i^2c^4,\notag\\
 u&=cT,& b_0&=c+R\Drat,\notag\\
 \Arat&=Q(\Drat u^2+b_0^2-y^2)+y^2-1,&
 \Brat&=2Qub_0,\notag\\
 P&=N_0N_1N_kN_t,&\RR&=\Arat^2-\Drat\Brat^2.
 \label{eq:rootdata}
\end{align}
There are exactly fourteen supplied positive unknowns in these
expressions. The two equations are
\begin{align}
 \boxed{\Phi_{14}=P(1+\RR)-1=0},&
 \qquad 75\text{ operations},\quad\deg\Phi_{14}=293,
 \label{eq:phi}\\
 \boxed{\Psi_{14}=[P(\Arat+1)-1]^2-\Drat(P\Brat)^2=0},&
 \qquad 73\text{ operations},\quad\deg\Psi_{14}=342.
 \label{eq:psi}
\end{align}
The shorter circuit for \(\Psi_{14}\) shares the products in the two
squared terms; evaluating the printed display naively need not attain
73. Both complete circuits are included.

\subsection{The signed-root issue must be resolved}
Algebraic elimination alone restores a nonzero integer \(f\) of either
sign. The existing signed-quotient theorem was written for the universal
compiler and its input bridge. It therefore needs an input-free
adaptation before it can be used here.

\begin{lemma}[Input-free signed quotient]\label{lem:signed}
On either valid coefficient tuple specified in this report, a zero of
\eqref{eq:normalized} with every supplied coordinate positive except
possibly \(T\in\Z\) necessarily has \(T>0\). Equivalently, allowing
only \(f\) to be signed while \(T>0\) also adds no zeros.
\end{lemma}
\begin{proof}
The unit exclusions and outer bounds in~\eqref{eq:outerbounds} do not
use \(T>0\). Keep \(C=0\) provisionally. Normalized strong divisibility
gives \(f=\chi_A(m)\), \(\psi_A(m)=ic^2\), and in fact
\(pc\mid m\): first \(p\mid m\) by strong divisibility, then
expansion of \((\chi_A(p)+c\sqrt\Delta)^{m/p}\) modulo \(c\)
gives \(c\mid m/p\).

Classify the \emph{absolute} auxiliary root rather than assume its
sign. The literal congruences \(V\equiv-c\pmod f\) and
\(V\equiv-R\pmod c\) remain valid. Squaring the first, and combining
it with the second and the bounds \(0<p,R<c/2\), gives \(p=R\)
as in the signed absorption proof~\cite{signedabs}. The no-wrap argument
needs only \(R+2<E\), which~\eqref{eq:outerbounds} supplies. Thus
\(2n=R+1\), both remaining units are positive, and transport gives
\(C>0\). In particular \(R\) is odd and \(C+F<q\).

Here are the power and population details that prevent a circular sign
argument. Put \(r=(R-1)/2\) and
\(\xi=(X+1)^{2r}/X^r\). Pell ratio bounds give
\[
 Y>\xi/2-1>X^r/3,\qquad a>X^{r+1}/3>2^R.
\]
The recurrence for \(\chi_A(j)-a\psi_A(j)\), whose initial values
are 1 and 2, gives \(\chi_A(j)-a\psi_A(j)\equiv2^j\pmod H\).
The main-root equation therefore gives \(X\equiv2^R\pmod H\).
Both positive numbers are below \(a<H\), so \(X=2^R\).
It follows that \(q=B^N\), with \(t=dN\), and \(q^3\mid X\).

The upper ratio estimate, using \(r\ge24\), is
\[
 0<c/k-\xi/2<16r/(X+1)<1/2.
\]
The exact expansion at \(X=2^{2r+1}\) is
\[
 \xi=M+\theta,\quad
 M=\binom{2r}{r}+\sum_{j=1}^r\binom{2r}{r+j}X^j,
 \quad0<\theta<1/4,
\]
where \(M\) is even. Since \(Y<c/k<Y+1\), we obtain \(Y=M/2\).
Thus \(\nu_2(Y)=\pc(r)-1\), and \(q^3\mid Y\) implies
\(\pc(R)=\pc(r)+1\ge3dN+2\). No sign of \(V\), or congruence
\(R\equiv3\pmod4\), has entered these deductions.

Apply the word identity~\eqref{eq:wordidentity} and the
inverse-population lemma. The masks~\eqref{eq:masksrecovered} follow.
Since \(M_C\equiv2\pmod4\) and \(J\equiv1\pmod4\), the low
mask forces \(C\equiv1\pmod4\). Reducing the source definition of
\(R\) gives \(R\equiv C+M_CJ\equiv3\pmod4\).

It remains to use the two unsquared congruences. Let
\(e=\operatorname{sign}(V)\) and classify the auxiliary Pell norm with
odd index \(\ell=2v+1\). Positive-sign \(\chi\) step-down gives
\(\ell=\epsilon R+2mz\), \(\epsilon\in\{1,-1\}\).
The two congruences imply, using \(c>2R\) and \(f>2c\),
\[
 e(-1)^{r+mz}=-1,\qquad e(-1)^{r+(m+1)z}=-1.
\]
Dividing proves \(z\) even; since \(R\equiv3\pmod4\), \(r\)
is odd and \(e=1\). Consequently
\(cTf=V+c+Rf^2>0\), so \(T>0\). This is the signed parity argument
of~\cite{signedsound}, with the old packed coordinate replaced exactly
by \(C\). The input bridge was not used.

Finally, the only direct occurrences of \(f\) in the parent are
\(f^2\) and \(Tf\). Simultaneously negating \(f,T\) preserves every
computed value. A negative \(f\) with positive \(T\) would contradict
the first assertion after this sign change. The strong equation excludes
\(f=0\).
\end{proof}

\subsection{Algebraic deletion and restoration}
Before imposing any equation, \(\Drat>c>0\). Since \(R\) is an
integer, \(b_0=c+R\Drat\ne0\); hence \(\Brat\ne0\).
At a positive zero of~\eqref{eq:phi}, the two integer factors are units.
Because \(\Brat\) is even,
\(1+\RR\equiv1+\Arat^2\pmod4\) cannot be \(-1\).
Therefore \(P=1\) and \(\RR=0\). The rational number
\(f=\Arat/\Brat\) satisfies \(f^2=\Drat\in\Z_{>0}\), so it is an
integer: a reduced rational square root of an integer has denominator one.
Then \(V=uf-b_0\) and
\[
 QV^2-(Q-1)y^2=\Arat+1-\Brat f=1,
 \qquad f^2-\Delta i^2c^4=1.
\]
Thus a signed normalized parent zero has been restored. Lemma~\ref{lem:signed}
makes \(f\) positive. Conversely, a positive normalized parent zero has
\(f^2=\Drat\), \(\Arat=\Brat f\), and \(P=1\), so deleting \(f\)
gives a zero of \(\Phi_{14}\). The restored root is unique.

For \(\Psi_{14}\), reduction of~\eqref{eq:psi} modulo \(P\) gives
\(P\mid1\), hence \(P=\pm1\). Put
\(f=(\Arat+1-P)/\Brat\). The equation again implies \(f^2=\Drat\)
and integrality. Its auxiliary norm is
\(QV^2-(Q-1)y^2=P\). As \(Q\) is a square, this is a square
modulo four and cannot be \(-1\). Thus \(P=1\), and the same
restoration applies. The two fourteen-variable polynomials have exactly
the same positive zero tuples. An exact all-value identity also holds:
\begin{equation}
 \Psi_{14}=P\Phi_{14}+(P-1)(2P\Arat-1).
\end{equation}
The equality of zero sets follows from the unit argument, not from this
identity alone.

\subsection{Uniform degrees and complete circuits}
Deleting the \(f\)-dependent cone from the 67-gate parent leaves exactly
50 live producers. The local norm finalizer adds 25 gates; the shared
outer norm adds 23. Thus
\[
 \Phi_{14}:44M+31A=75,\qquad\Psi_{14}:44M+29A=73.
\]
The factored computation uses
\(\Drat\Brat^2=4(Q\Drat u^2)(Qb_0^2)\); there is no division gate.
With \(q_0,k_0,T_0\) as above, the exact leading forms are
\begin{align}
 (\Phi_{14})_{293}
 &=8\gamma h i^{12}k_0^{27}w^{21}s^{48}
   q_0^{177}(q_0-F)^4T_0,\notag\\
 (\Psi_{14})_{342}
 &=64\gamma^2h^2i^{12}k_0^{30}w^{26}s^{56}
   q_0^{206}(q_0-F)^4T_0^2.\label{eq:rootleaders}
\end{align}
The respective distinguished monomials and their coefficients are
\[
\begin{array}{c|c}
 J^{182}\gamma h i^{12}\eta^{27}w^{21}s^{48}z&-8\beta^{182}\\
 J^{212}\gamma^2h^2i^{12}\eta^{30}w^{26}s^{56}z^2&64\beta^{212}.
\end{array}
\]
Both coefficients are uniformly nonzero.
An independent sparse-polynomial audit reconstructs these complete
highest homogeneous forms exactly.

\section{Twenty-four quadratics in thirty-three unknowns}\label{sec:quadratic}
Start with the fifteen unknowns of Section~\ref{sec:system}, and add
the eighteen positive registers
\[
 N,E,P_*,c,a,X,D,C_2,G,L,R,F_2,V_0,V,g,I,B_a,\Delta.
\]
In this section use only the affine abbreviations
\(q=\beta J+1\), \(k=\eta+\zeta\), \(H=4a+3\),
and \(C=q-F-\alpha\). The complete quadratic system is:
\begin{center}\small
\begin{longtable}{rl@{\qquad}rl}
\toprule
1.&\(L=q^2-1\)&2.&\(N=q(L+1)\)\\
3.&\(X=wq\)&4.&\(a=E+sN\)\\
5.&\(E=X(a-E)\)&6.&\(c=k(a-E)+\eta\)\\
7.&\(P_*=E(c-\eta)+k\)&8.&\(\tau^2=P_*(P_*-k)+1\)\\
9.&\(D=X+ac+\gamma H\)&10.&\(\Delta=a^2+H\)\\
11.&\(C_2=c^2\)&12.&\(D^2-\Delta C_2=1\)\\
13.&\(k=R+1+hE\)&14.&\(G=(q-1)(q-F)+\alpha\)\\
15.&\(R=LG+(M_C+q\Msrc)J\)&&\\
16.&\multicolumn{3}{l}{\((K+w)C+(q-F)-z(q-1)=1\)}\\
17.&\(F_2=f^2\)&18.&\(V_0=Tf-1\)\\
19.&\(V=cV_0-RF_2\)&20.&\(g=y^2-V^2\)\\
21.&\(I=iC_2\)&22.&\(B_a=I^2\)\\
23.&\(B_a=\Delta(F_2-1)\)&24.&\(y^2-B_ag=1\)\\
\bottomrule
\end{longtable}
\end{center}
Let \(r_j\) be left side minus right side. Define the literal polynomial
\begin{equation}
 \boxed{Q_{33}=\sum_{j=1}^{24}r_j^2.}\label{eq:q33}
\end{equation}
The sum has a fixed 24 terms. It contains no hidden variable-length
conjunction or computation predicate. The coefficient of \(\tau^4\)
is 1, so its degree is exactly four.

The projection follows by successive substitution: \(N=q^3\),
\(Y=a-E=sq^3\), \(E=XY\), and \(P_*=EkY+k\). The six
remaining constraints are precisely~\eqref{eq:system}. Conversely
every source solution uniquely determines all eighteen registers.
The less immediate positivity claims follow from the preceding proof:
transport gives \(C>0\), hence \(G,R>0\); the strong-rank argument
gives \(f>2c\), hence \(V_0>0\); auxiliary recovery gives \(V>1\)
and \((B_a-1)g=V^2-1>0\), hence \(g>0\).

The saving over package 08 is exact. Its scaled-root equation used
\(I=i\Delta c^2\) and a register \(K_2=D^2-1=\Delta c^2\).
After using the ordinary root \(I=ic^2\), the register \(K_2\) is
needed only by the main norm. Combining its two equations into
\(D^2-\Delta C_2=1\) removes one register and one equation.

\section{Domains: nonnegative and unrestricted integers}\label{sec:domains}
All principal equations have been stated over \(\Pos\). Interpreting
them unchanged over \(\Z\) would be incorrect. For a nonnegative
domain, replace every positive unknown \(x\) by \(n_x+1\);
this preserves the variable count and degree, though a literal evaluator
must charge the extra additions.

For unrestricted integers use
\begin{equation}
 \operatorname{Pos}(u,v,t)=1+u^2+v^2+t(t+1).\label{eq:pos}
\end{equation}
It is always positive and represents every positive integer. Indeed
\(4n-3\equiv1\pmod4\) is represented by three squares by the
three-square theorem; exactly one square is odd. Writing it as
\((2t+1)^2\) and the other two as \((2u)^2,(2v)^2\) gives
\eqref{eq:pos}. This also covers \(n=1\).

In both fourteen-variable polynomials, \(\tau,y\) occur only squared,
and neither can vanish at a zero. The first norm excludes \(\tau=0\);
the restored auxiliary equation with \(Q>1\) excludes \(y=0\).
Thus these two may already range over all integers. Substitute
\eqref{eq:pos} for each of the other twelve coordinates. The resulting
single equations have 38 unrestricted integer unknowns:
\begin{center}\small
\begin{tabular}{lrrr}
\toprule
Integer polynomial & Unknowns & Exact degree & Operations\\
\midrule
\(\Phi_{14}\) after substitution & 38 & 586 & 159\\
\(\Psi_{14}\) after substitution & 38 & 684 & 157\\
\bottomrule
\end{tabular}
\end{center}
Each positivity form costs seven gates. The leading forms in
\eqref{eq:rootleaders} contain neither \(\tau\) nor \(y\), so the
degree doubles exactly. Substitution by disjoint positive quadratic
forms preserves a nonzero leading form, as seen by setting two of each
triple to zero and obtaining the injective monomial substitution
\(x\mapsto u_x^2\).

There are also lower-degree integer alternatives. In the quadratic lift,
only thirteen original coordinates need positivity; \(\tau,y\) and
all eighteen registers may be signed. The equations determine the
registers and the preceding positivity proof makes them positive at
every solution. Retain the 33 variables over \(\Z\), and introduce
three new variables and one equation \(x=\operatorname{Pos}(u,v,t)\)
for each of those thirteen coordinates. This gives 37 quadratics in
72 unrestricted integers, or one quartic sum of squares. Direct
substitution instead removes the thirteen old coordinates and introduces
39, giving 59 integer variables and a degree-eight sum of squares.
These improve particular explicit integer lifts in the eight packages;
they do not improve the classical universal variable/degree frontier.

\section{Smaller fixed coefficients from four directed cells}\label{sec:coefficients}
This section specifies all four numerals used above. The complete finite
predicate and exact recipes are also executable in
\file{code/four_cell_compiler.py}; the scalar values are recorded in
\file{data/four_cell_constants.json}. No part of the definition asks
whether K215 halts.

\subsection{The exact machine and the local alphabet}
Use the 430-row table \file{vendor/k215.tm}, with 215 working states,
two symbols, initial state 0 and a two-way blank tape. Its SHA-256 is
\begin{center}\small
\texttt{fdbcef0825886035c3091c8537a69535bc96f9e265f0ae45863543d40ae70090}.
\end{center}
The halt state is 215; for tableau padding add the two stationary
identity actions at that state. The actual first action is
\((0,0)\mapsto(1,\mathrm R,93)\).

The cell alphabet has \(a_0=439\) symbols. Codes 0,1 are headless
tape bits. Code \(2(q+1)+b\) is bit \(b\) with a head in state
\(q\in\{0,\ldots,215\}\). Codes 434,435,436 are initial phases
\(L,Q,R\), with blank payloads and an initial-state head only at
\(Q\). Codes 437,438 are vertical and horizontal boundaries \(V,H\).
The old local predicate inspected a nine-cell square but depended only
on the five-cell cross. Its preceding-time cell was used solely to
detect an \(H\) immediately before an initial phase.

Move that test to its predecessor. A new tile is
\((\mathrm{left},\mathrm{center},\mathrm{right},\mathrm{following})\).
At an \(H\) center require a following phase \(L,Q,R\). At an
interior center forbid an initial phase on a following interior cell.
The horizontal initial grammar remains
\(L\to L\text{ or }Q\), \(Q\to R\), \(R\to R\), with
\(V\) forcing \(L\) at the left edge and \(R\) at the right edge.
Vertical boundaries propagate vertically. The horizontal boundary bit
propagates across each nonvertical strip; consecutive horizontal
boundaries and consecutive vertical boundaries are forbidden.

For an interior-to-interior update, write the current head's prescribed
bit, if any, and collect incoming heads from a left neighbor moving
right, a stationary center head, or a right neighbor moving left. Two
or more incoming heads are forbidden, and the following payload must
equal this update. A vertical-boundary neighbor supplies blank headless
payload; a horizontal-boundary neighbor is excluded. A head cannot
cross a vertical boundary. A following \(H\) is allowed only when the
center head, if present, is halted; this final boundary step does not
perform an interior update. These are the literal tests in the provided
predicate.

\begin{lemma}[Four-cell equivalence and reconstruction]\label{lem:fourcells}
The old and new local predicates define the same periodic cell planes.
For neighboring tiles \(C,R_t,Y_t\), the three equalities
\[
 \mathrm{right}(C)=\mathrm{center}(R_t),\qquad
 \mathrm{center}(C)=\mathrm{left}(R_t),\qquad
 \mathrm{following}(C)=\mathrm{center}(Y_t)
\]
are sufficient to reconstruct such a plane from a cyclic tile word.
\end{lemma}
\begin{proof}
For any interior cell, its predecessor is \(H\) exactly when the new
rules force that cell to have an initial phase: an interior predecessor
forbids such a phase, and a \(V\) predecessor forces \(V\) rather
than an interior cell. All other old clauses are unchanged. This
reindexing is valid even if offsets in a cyclic presentation coincide.

For a cyclic word \(T_i\) with neighbor offsets \(-1,-h_0\), put
\(u_i=\mathrm{center}(T_i)\). The equalities give the other entries
as \(u_{i+1},u_{i-1},u_{i-h_0}\), respectively. Pull back along
\((x,t)\mapsto-x-h_0t\pmod N\) to obtain a periodic plane. No
coprimality assumption is needed. Genuine neighborhoods give the
converse.
\end{proof}

The distinguished first tile is
\(S=(434,435,436,1)\), recording the literal first transition.
The distinguished second tile is \(E=(437,437,0,437)\). It is
valid at its center but has two adjacent vertical boundaries, so it has
no global extension. Enumerate these two first and then all other
locally valid four-tuples in lexicographic order. The exact number is
\begin{equation}
 \kappa=71\,356\,645.\label{eq:tilecount}
\end{equation}
The transition table has 348 right-moving and 82 left-moving actions.
Adding the initial phase gives incoming-head counts 349 and 82 and
nonincoming nonhorizontal counts 89 and 356. The center-class census is
\begin{center}\small
\begin{tabular}{lrr}
\toprule
Center class & Multiplicity & Valid tiles per center\\
\midrule
Vertical boundary & 1 & 192282\\
Horizontal boundary & 1 & 2634\\
Headless bit & 2 & 355508\\
Halt head & 2 & 223966\\
Working head & 430 & 162788\\
Initial \(L\) & 1 & 1754\\
Initial \(Q\) & 1 & 437\\
Initial \(R\) & 1 & 1750\\
\bottomrule
\end{tabular}
\end{center}
The accompanying proof and checker refine the working-head row into
four direction/write classes and verify every left/right pair for each
representative class. The predecessor multiplicity and the unused
four-corner multiplicity have both disappeared.

A marked periodic plane still forces a proper finite strip containing
the initial \(L^uQR^v\) row and exactly one initial head in that strip.
Temporal recurrence supplies a later boundary; the last interior head
must be halted. Conversely any finite halting run admits independent
width and height padding and a helical cyclic presentation. It has one
Start and no forbidden tile. A zero temporal stride is impossible since
the Start center \(Q\) differs from its following center 1. Thus
Lemma~\ref{lem:fourcells} preserves both directions of the blank-machine
tableau theorem.

\subsection{A finite recipe for every coefficient}
The geometry below uses symbols distinct from the Pell unknowns. Let
\(\ell=27\), \(A_0=2^\ell\), and take the ordered slot list
\[
 \mathcal L=((1,0),(1,1),(1,2),(2,1)).
\]
For slot \((r,c)\) and symbol \(v<a_0\), put
\(e(r,c,v)=\kappa+3a_0+(2-r)3a_0+(2-c)a_0+v\).
Define
\begin{align*}
 p_0&=35676757,&j_{\max}&=4a_0+4+p_0=35678517,\\
 \mu&=A_0-2+\sum_{j=1}^{j_{\max}}A_0^j,&
 m_0&=2\ell+11a_0+8+2p_0=71358405,\\
 e_*&=\kappa+12a_0=71361913,& M_0&=\kappa+15a_0+1=71363231.
\end{align*}
The native support \(\mathcal E\) is the disjoint union of selectors
\(0,\ldots,\kappa-1\), the \(4a_0\) payload positions, the single
dummy position \(e_*\), and anchors \(M_0,3M_0,9M_0,27M_0\).
It has \(m_0+1\) elements. For tile \(T_i\), define
\begin{align*}
 c_i&=1+\sum_{j=0}^3 A_0^{1+ja_0+T_i[j]}
              +\sum_{t=1}^4A_0^{4a_0+t},\\
 c_{e(\mathcal L_j,v)}&=A_0^{1+ja_0+v},\qquad
 c_{3^{t-1}M_0}=A_0^{4a_0+t},\qquad c_{e_*}=0.
\end{align*}
The remaining fixed integers are
\begin{align*}
 E_{\max}&=27M_0=1926807237,&
 H_C&=51M_0+3a_0+1=3639526099,\\
 T_1&=H_C+2E_{\max}+a_0+1=7493141013,&
 T_2&=T_1+4E_{\max}+2=15200369963,\\
 L_C&=5^{15}=30517578125,& b&=5^{13}=1220703125,\\
 \rho&=2^b,&d&=bL_C=5^{28},\qquad B=\rho^{L_C}=2^d.
\end{align*}
All sums below are over fixed finite sets. Define
\begin{align}
 D_C={}&\rho^{3a_0}+\rho^{H_C+a_0}+\rho^{8M_0}+\rho^{24M_0}
       +\sum_{e\in\mathcal E}c_e(\rho^{T_1-e}+\rho^{T_2-e}),\notag\\
 D_R={}&\rho^{H_C},\notag\\
 M_F^0={}&\mu(\rho^{T_1}+\rho^{T_2})
       +\sum_{v<a_0}\rho^{e(1,1,v)}\notag\\
 &+\sum_{v<a_0}\bigl(\rho^{H_C+e(1,0,v)}+
                         \rho^{H_C+e(1,1,v)}\bigr)
       +\rho^{9M_0}+\rho^{27M_0},\notag\\
 M_C={}&B-1-\sum_{\substack{e\in\mathcal E\\e\ne1}}\rho^e
                 -2\rho^{e_*},\qquad \Mnat=M_F^0+4.
\end{align}
The four coefficients of every principal construction are exactly
\begin{equation}
 \boxed{\beta=B-1,\qquad K=D_C+BD_R,\qquad M_C,
                  \qquad\Msrc=\Mnat+B-1.}\label{eq:ports}
\end{equation}
The transition table and finite lexicographic enumeration make these
specific integers. No free program index, unspecified existential
coefficient, or oracle occurs in~\eqref{eq:ports}.

\subsection{Why the smaller layout satisfies the compiler hypotheses}
The selector clause uses only \(\kappa<A_0\), a common low selector
coefficient 1, and disjoint powers of \(A_0\) for payloads and anchors.
Its proof is independent of whether there are nine or four payload
positions. The vertical target \(e(1,1,v)\) receives the following
payload through shift \(3a_0\); the temporal word supplies its center.
The horizontal targets at \(H_C+e(1,0,v)\) and \(H_C+e(1,1,v)\)
give the two horizontal equalities. The anchor targets remain
\(9M_0,27M_0\). These are precisely the three copy conditions of
Lemma~\ref{lem:fourcells} and the original alignment conditions.

All nonanchor native support lies below \(M_0\). The copy bands and
the two center bands are separated by the displayed inequalities
\[
 T_2-T_1>2E_{\max},\quad
 L_C-(T_2-T_1)>E_{\max},\quad T_2+E_{\max}<L_C.
\]
There are no accidental contributions to the listed targets. The anchor
difference \(18M_0\) occurs in the native support in the unique way
\(27M_0-9M_0\); it still forces the temporal rotation to align to a
whole cell. An occupied cell cannot agree with an empty neighbor at its
one-hot copy coordinates, so occupancy propagates from Start around the
entire cyclic word. No diagonal copy condition is needed.

The fixed value \(b=5^{13}\) exceeds the sufficient radix threshold
963319989. In particular
\[
 \rho\ge\max\{4(m_0+3)(2\sum_e c_e+6)+8,\ 2\mu+4,\ 16\}.
\]
This is the same conservative no-carry bound as in the parent compiler.
The extra dummy bit's spill stays below \(M_0\); at least one of the
two center bands is clean under any initial binary rotation. The two
anchor parity tests then exclude every nonzero within-radix rotation,
and their unique difference excludes every nonzero within-cell rotation.
The inherited clean-band and anchor proof therefore applies literally
to the new support~\cite{blank,compiler}.

Population accounting changes from twelve to three copy blocks:
\[
 \pc(\mu)=\ell+4a_0+3+p_0,\quad
 \pc(M_F^0)=2\pc(\mu)+3a_0+2=m_0.
\]
Consequently
\begin{gather}
 0<M_C,\Mnat<B-1,\quad M_C\equiv2\pmod4,
 \quad\Mnat\equiv4\pmod8,\notag\\
 \pc(M_C)=d-m_0-1,
 \quad\pc(\Mnat)=m_0+1,
 \quad\pc(M_C)+\pc(\Mnat)=d.
\end{gather}
These are the exact mask conditions used in the arithmetic and signed
proofs. Since \(T_2-T_1\equiv2\pmod4\), \(a_0\equiv3\pmod4\)
and \(\rho\equiv2\pmod5\), the paired center terms cancel modulo
five and
\begin{equation}
 2D_C-D_R\equiv3\pmod5.\label{eq:fiveunit}
\end{equation}
Thus no conditional correction coefficient is needed.

Completeness must tune the actual packed index, rather than choose an
independent exponent. Take independently padded powers-of-five width
and height with \(N>25d\), and temporal stride \(h_0\). Changing a
low dummy bit in cell \(j\) changes the index by \(-\Gamma B^j\),
where
\[
 \Gamma=\rho^{e_*}(q^2-1)
       [1+q(D_C+BD_R+B^{h_0})]
\]
is invertible modulo \(dN\) by~\eqref{eq:fiveunit}. The inherited
Boolean subset lemma tunes \(R\equiv dh_0\pmod{dN}\) without
altering any tile or copy clause. The canonical no-carry margin still
gives \(2C+F<q/3\), so \(\alpha=q-F-C>0\).
The normalized Pell converse supplies all remaining positive witnesses.
This establishes the new coefficient compiler in both directions.

\subsection{What the height improvement does and does not say}
The old raw-K215 coefficient \(B-1\) had \(5^{67}\) binary digits.
The new coefficient has \(5^{28}\), so its bit length is smaller by
the exact factor
\[
 5^{39}=1\,818\,989\,403\,545\,856\,475\,830\,078\,125.
\]
The exact new bit lengths are
\begin{center}
\begin{tabular}{lr}
\toprule
Coefficient & Binary digits\\
\midrule
\(\beta\) and \(M_C\) & \(37\,252\,902\,984\,619\,140\,625\)\\
\(\Msrc\) & \(37\,252\,902\,984\,619\,140\,626\)\\
\(K\) & \(41\,695\,683\,867\,187\,500\,001\)\\
\bottomrule
\end{tabular}
\end{center}
For \(K\), use \(D_C<B\) and
\(K=2^{b(L_C+H_C)}+D_C\). The coefficients remain far too large to
expand in the archive. This is an exact finite specification with a
substantial height reduction, not a claim of small decimal coefficients
or an efficiently materializable polynomial.

\section{Logical and executable verification status}\label{sec:status}
The source snapshot is pinned on the title page. The repository records
closed Lean and Rocq endpoint checks for the literal K215 table; the
Lean statement is
\file{KEndpoint215.blank_halts_iff_zfc_refutation}. We inspected the
statement and the repository's verification record. Those proof-assistant
builds were not rerun. The theorem refers to the effective predicate
\file{BProv}; ordinary effective proof translations connect it to the
chosen conventional formulation of \(\Con(\mathrm{ZFC})\).

The compiler and all improvements in this report have ordinary
mathematical proofs, not new kernel-checkable proofs. Their arguments
use finite strings and tables, integer arithmetic, Pell recurrences,
finite combinatorics, and standard number theory. The K215 calibration
is supplied by the repository's ordinary proof which the formal endpoint
mirrors. These arguments can be carried out in ZFC. No appeal to
``Lean is sound, therefore ZFC proves the statement'' is made.

The included verifiers do establish concrete, reproducible facts:
\begin{itemize}
\item exact binding of all 430 transition rows and the dense-table hash;
\item the four-cell census, scalar layout and mask populations;
\item local/copy checks on 32 padded halting-control tableaux, rejection
of a truncated nonhalting prefix, and 10,566 sparse toy support checks;
\item gate topology, liveness, complete operation ledgers, and literal
agreement of the arithmetic packets with their formulas;
\item exact all-ring root-elimination identities and exact independent
sparse proofs of the degree-293 and degree-342 leading forms;
\item all 24 quadratic residuals, their symbolic projection identities,
the register counts, and the stated integer-domain conversions.
\end{itemize}
They do not decide whether K215 halts, materialize the four giant
coefficients, or prove the compiler theorem by finite testing. The small
coefficient substitutions used in arithmetic diagnostics are not claimed
to encode K215 or any accepting computation. Halting-control machines in
the local tableau checks are deliberately modified test machines.

If ZFC is consistent, the equivalences do not allow it to prove these
equations have no solutions: that would give its own consistency.
Mere consistency alone is not asserted to give the stronger two-sided
unprovability of both solvability and unsolvability. The result is a
representation of the consistency assertion, not a consistency proof.

\section{Further research questions}
\begin{enumerate}
\item \textbf{Formalize a reusable signed quotient theorem.} The delicate
steps are the initially possible \(C=0\) boundary, power recovery before
parity, and use of both unsquared auxiliary congruences. A theorem stated
at this interface would cover the input-free and universal cases.
\item \textbf{Formalize the ordinary strong-rank interface.} Its conclusion
\(p\mid m\), not only \(c\mid m\), is what restores the deleted
quotient integrally. Packaging both conclusions prevents a recurring
proof-order gap.
\item \textbf{Compress the local alphabet further.} Mark ordinary cells as
left or right of the unique head and enforce a local
\(V\,L^*\,\mathrm{Head}\,R^*\,V\) grammar. This may reduce legal
neighborhoods to a much smaller family at the price of a few extra
symbols. No such compiler is claimed here.
\item \textbf{Use the C222 proof calculus directly.} A direct local
calculus or stack encoding might reduce coefficient height much more
than shaving states from the already small machine. It needs a complete
fixed-arity encoding and a calibrated consistency equivalence.
\item \textbf{Reduce fourteen unknowns without a degree explosion.} The
present elimination removes one root but raises degree. Can an ordinary
strong-norm elimination preserve its lower degree while controlling
rational-root denominators and signs?
\item \textbf{Reduce the 33-register lift.} Search for another quadratic
relation allowing an intermediate register to be removed without adding
an inequality or silently imposing positivity.
\item \textbf{Make a materializable coefficient benchmark.} Operations,
degree, and variable count omit the dominant storage cost. Compare
constructions with explicit bounds on coefficient bit length and on the
time needed to produce their expanded coefficients.
\item \textbf{Extract a concrete nine-unknown specialization.} The abstract
existence result is already known; an explicit specialization with a
fully bound index and a reproducible coefficient description would allow
an honest comparison with the constructions here.
\end{enumerate}

\appendix
\section{Files, reproduction, and provenance}
The archive is a self-contained integration directory. It includes the
LaTeX source and compiled PDF; the exact K215 table; the original
67-operation circuit as immutable JSON data; every newly emitted circuit;
the expanded quadratic systems; detailed proof notes and independent
reviews; selected unchanged upstream reference notes; and standard-library
Python verifiers. The scripts read the vendored parent data rather than
executing inherited repository programs. The directory is self-contained
for these finite checks and the PDF build, but does not include the full
upstream Lean/Rocq development or every historical proof dependency.
\begin{lstlisting}
python3 verify.py
pdflatex -interaction=nonstopmode -halt-on-error report.tex
pdflatex -interaction=nonstopmode -halt-on-error report.tex
\end{lstlisting}
Individual builders compare the committed data by default. Their
\file{--write} options regenerate only their own packets. All new fixed
coefficients are the tuple in \file{data/four_cell_constants.json}.
The manifest records file hashes and the source commit. Supplied license
and attribution notices preserve the original machine contributors.

The ordinary system and quadratic variants build on packages 02, 06,
and 08. The fourteen-variable construction builds on package 06, package
05, and the repository's complete90/92 norm constructions. The four-cell
recipe modifies package 06's local compiler while preserving its
coefficient conventions and arithmetic interface. The improvements are
not attributed to the inherited source authors; their prior contributions
are identified separately throughout.

\begin{thebibliography}{99}\small\raggedright
\bibitem{reconcile} ProveIt, \emph{Con(ZFC) in Diophantine form},
reconciliation of eight packages, pinned snapshot above.
\url{https://github.com/VladimirReshetnikov/ProveIt/tree/3b201cdea32d2c816f5f3d20734051f25fd3ef14/SetTheory/Cardinals/docs/reports/hilbert-tenth-problem/con-zfc-diophantine-forms}.
\bibitem{endpoint} ProveIt, \emph{Formal correctness of the exact
215-state machine}; and \file{Computability/BusyBeaver/Lean/BusyBeaver/ZFC/KEndpoint215.lean}.
\url{https://github.com/VladimirReshetnikov/ProveIt/tree/3b201cdea32d2c816f5f3d20734051f25fd3ef14/Computability/BusyBeaver/Research/ZFC/k225-formal/verification/k215}.
\bibitem{blank} Package 06, \emph{Explicit Diophantine equations for
Con(ZFC)}, \file{06-blank67/report.tex}, and
\file{06-blank67/proofs/fixed_blank_67_proof.md}, in~\cite{reconcile}.
\bibitem{ordinary} Package 02, \file{PELL_INTERFACE_LEMMAS.md} and
\file{FIXED_INPUT75_ORDINARY_PROOF.md}, in~\cite{reconcile}.
\bibitem{fixed15} Package 05, \emph{Fixed-input Diophantine forms of
ZFC consistency}, \file{05-fixed15-16/paper/fixed_input_zfc.tex},
in~\cite{reconcile}.
\bibitem{quartic34} Package 08,
\file{evidence/diophantine_zfc_encoding/start_only_quartic34_proof.md},
in~\cite{reconcile}.
\bibitem{relaxed} ProveIt, \emph{Recovering integrality after relaxing an
auxiliary Pell norm},
\file{Computability/HilbertTenthProblem/Papers/1980/PELL_RELAXED_AUXILIARY_PROOF.md}.
\bibitem{compiler} ProveIt, native stream/queue compiler research notes,
especially \file{complete75_half_binomial_compiler.md} and the
asymmetric-scale and normalized strong-output proofs.
\url{https://github.com/VladimirReshetnikov/ProveIt/tree/3b201cdea32d2c816f5f3d20734051f25fd3ef14/Computability/HilbertTenthProblem/Papers/research-wip/native-stream-queue}.
\bibitem{signedabs} ProveIt,
\file{complete84_signed_quotient_absorption.md}, in the directory of~\cite{compiler}.
\bibitem{signedsound} ProveIt,
\file{complete84_signed_quotient_soundness_scout.md} and
\file{review_complete84_signed_quotient_soundness.md}, in~\cite{compiler}.
\bibitem{norm90} ProveIt,
\file{complete90_signed_root_elimination.md}, in~\cite{compiler}.
\bibitem{norm92} ProveIt,
\file{complete92_norm_unit_guard.md}, in~\cite{compiler}.
\bibitem{nine} ProveIt, \file{Jones1982.nine_unknowns},
\file{Computability/HilbertTenthProblem/Lean/Diophantine/Paper1982/NineMain.lean}.
The cited repository theorem is an existence result; no concrete
nine-variable equation is emitted by this report.
\end{thebibliography}
\end{document}
